% =====================================================================
% v26 §III extension: Theoretical Foundation of AI Necessity
% Theorem 8 = Main result (AI Necessity)
% Builds: Theorem 4 (Static Regret LB), Theorem 5 (Sufficient Statistic Gap),
%         Theorem 6 (Memory Necessity), Theorem 7 (Analytical Intractability)
% =====================================================================

\section{Theoretical Foundation: Why ML Is Logically Necessary}
\label{sec:theoretical-foundation}

The Linear-Score Ceiling Theorem (Theorem~\ref{thm:ceiling}) establishes
that linear classifiers on $(CC, RTT)$ are reduced to chance under
moment-matched adversaries. This is a strong claim about one detector
class. Skeptical reviewers correctly ask: is the result narrow to a
particular adversary, or does it generalise to other operational regimes?

This section places Theorem~\ref{thm:ceiling} inside a broader
control-theoretic and information-theoretic framework. We prove four
additional theorems (\ref{thm:static-regret}--\ref{thm:intractability})
that together establish a main result, the \emph{AI Necessity Theorem}
(Theorem~\ref{thm:ai-necessity}): under three operationally realistic
assumptions, no policy in the constant-parameter, linear, or memoryless
function classes can achieve bounded regret against the optimal
sub-leader controller, while machine learning is the only known
function class with sample-complexity-bounded guarantees.

\subsection{Formal Model}

Consider an $N$-node consensus cluster operating at heartbeat ticks
$t \in \mathbb{N}$. At each tick:
\begin{itemize}\setlength\itemsep{1pt}
\item State $s_t \in \mathcal{S}$ collects per-node measurements
$(CC_i, RTT_i, \text{logTerm}_i, \text{logIdx}_i, \ldots)$.
\item Disturbance $w_t \sim P_t$, where $P_t$ may be \emph{time-varying},
encodes network jitter, GC stalls, traffic shifts, and Byzantine behavior.
\item Control input $u_t = \pi(s_{1:t}) \in \mathcal{U}$ is the
sub-leader designation produced by policy $\pi$.
\item Cost $c_t = \mathbb{1}[\text{deadline miss in tick } t] +
\lambda \cdot \mathbb{1}[\text{cascade event in tick } t]$, with
$\lambda \gg 1$ (we use $\lambda = 100$).
\end{itemize}
The long-run-average objective is
\begin{equation}
  \label{eq:objective}
  J(\pi) = \limsup_{T \to \infty} \frac{1}{T} \mathbb{E}\!\left[ \sum_{t=1}^T c_t \;\Big|\; \pi \right].
\end{equation}
Let $\pi^*$ denote the Bayes-optimal policy with full knowledge of $P_t$.

\subsection{Policy Classes}

We distinguish:
\begin{itemize}\setlength\itemsep{1pt}
\item $\Pi_{const}$: constant policies $\pi(s) = u^*$.
\item $\Pi_{lin}^{mem-0}$: linear scoring on instantaneous state with
fixed weights, $\pi(s) = \arg\max_i (w_{cc} CC_i(t) + w_{rtt} RTT_i(t))$.
\item $\Pi_{lin}^{mem-K}$: linear scoring on window-$K$ statistics
(e.g., S-Raft's $W{=}100$ windowed $CC$).
\item $\Pi_{poly}^{mem-K}$: degree-$d$ polynomial of windowed statistics.
\item $\Pi_{NN}$: neural-network policies (e.g., the Transformer of
Section~\ref{sec:predictor}).
\item $\Pi^*$: the Bayes-optimal policy.
\end{itemize}
We have the strict inclusion
$\Pi_{const} \subsetneq \Pi_{lin}^{mem-0} \subsetneq \Pi_{lin}^{mem-K}
\subsetneq \Pi_{poly}^{mem-K} \subsetneq \Pi_{NN}$, with the question of
whether $\Pi_{NN}$ approaches $\Pi^*$ being precisely what learning
theory studies.

\subsection{Theorem 4: Static Regret Lower Bound}

\begin{theorem}[Static Regret under Switching Disturbance]
\label{thm:static-regret}
Let the disturbance distribution be a switching process
$P_t \in \{P_A, P_B\}$ with mode-switch probability $\rho > 0$ per tick.
Let $w^{(M)} = (w_{cc}^{(M)}, w_{rtt}^{(M)})$ be the Bayes-optimal score
weights for mode $M \in \{A, B\}$. If $\|w^{(A)} - w^{(B)}\|_2 \ge
\epsilon > 0$, then for every constant-weight policy
$\pi \in \Pi_{lin}^{mem-K}$:
\begin{equation}
  J(\pi) - J(\pi^*) \;\ge\; \frac{\epsilon^2}{8 L^2} \cdot \rho,
\end{equation}
where $L$ is the Lipschitz constant of the score function.
\end{theorem}

\begin{proof}[Proof sketch]
Any constant weight $w^\pi$ has $\|w^\pi - w^{(M)}\|_2 \ge \epsilon/2$ for
at least one mode $M$. In that mode, the score-induced classifier
operates at a sub-optimal margin, yielding excess cost
$\ge \|w^\pi - w^{(M)}\|^2 / (2 L^2)$. The cluster spends fraction
$\ge \rho$ of its time in that mode in steady state.
\end{proof}

\textbf{Interpretation.} No single hand-tuned weight setting can be
optimal for both regimes simultaneously, and any real operation has
multiple regimes (diurnal, GC, deploy rollouts, attack vs no-attack).
The phrase ``just tune the weights better'' is logically refuted by
Theorem~\ref{thm:static-regret}.

\subsection{Theorem 5: Sufficient Statistic Gap}

\begin{theorem}[Information Capacity of Linear Detectors]
\label{thm:suff-stat}
Let $P_L$ and $P_B$ be the joint distributions of measurement vector
$X \in \mathbb{R}^d$ under legitimate and Byzantine nodes, with matched
first- and second-order moments but distinct higher-order cumulants
(e.g., distinct fourth cumulants). Then for any linear function
$f(X) = w^\top X + c$:
\begin{equation}
  I\!\left(f(X);\, \mathbb{1}[X \sim P_B]\right) = 0,
\end{equation}
while there exists a non-linear function $g : \mathbb{R}^d \to \mathbb{R}$
with $I(g(X); \mathbb{1}[X \sim P_B]) > 0$.
\end{theorem}

\begin{proof}
The marginal of $f(X)$ is determined by its mean and variance, which
agree across $P_L$ and $P_B$ by hypothesis. The pushforward measures
coincide, hence the mutual information vanishes. Conversely, for
$g(X) = X_1^4$ the marginal is determined by the fourth moment, which
differs between $P_L$ and $P_B$, giving positive mutual information.
\end{proof}

\textbf{Interpretation.} Theorem~\ref{thm:suff-stat} is the
information-theoretic dual of Theorem~\ref{thm:ceiling}: the linear
detector carries \emph{zero bits} of Byzantine information against a
moment-matched attacker. No threshold tuning recovers what was never
encoded.

\subsection{Theorem 6: Memory Necessity from Autocorrelation}

\begin{theorem}[Memoryless Policies Lose Autocorrelation Signal]
\label{thm:memory}
Let $\{X(t)\}$ be a stationary process with $\mathrm{AR}(1)$
autocorrelation $\rho_{AR} \neq 0$ for legitimate nodes, while
$\{X(t)\}$ is i.i.d.\ ($\rho_{AR} = 0$) for Byzantine nodes. For any
memoryless function $f : \mathcal{X} \to \mathbb{R}$:
\begin{equation}
  I\!\left(f(X(t));\, \mathbb{1}_B\right) \le I\!\left(X(t);\, \mathbb{1}_B\right),
\end{equation}
while for any windowed function
$g : \mathcal{X}^K \to \mathbb{R}$ ($K \ge 2$):
\begin{equation}
  I\!\left(g(X(t-K+1), \ldots, X(t));\, \mathbb{1}_B\right)
  \;>\; I\!\left(X(t);\, \mathbb{1}_B\right).
\end{equation}
The inequality is strict whenever $\rho_{AR} \neq 0$.
\end{theorem}

\begin{proof}[Proof sketch]
Autocorrelation $\rho_{AR}$ is invisible in the marginal distribution
of $X(t)$ but visible in the lag-1 covariance. By the data-processing
inequality, any memoryless function cannot exceed the marginal
information. A windowed function admits the lag-1 covariance as a
feature, strictly increasing the available information whenever
$\rho_{AR} \neq 0$.
\end{proof}

\textbf{Interpretation.} S-Raft's $W{=}100$ window already activates the
memory requirement, but the linear combination on top of the window
still suffers Theorem~\ref{thm:suff-stat}'s loss. The combination of
\emph{memory} and \emph{non-linearity} is required.

\subsection{Theorem 7: Analytical Intractability}

\begin{theorem}[The Optimal Policy Has No Closed Form]
\label{thm:intractability}
Under the assumptions of
Theorems~\ref{thm:static-regret}--\ref{thm:memory}, the Bayes-optimal
policy $\pi^*$ depends on the latent parameters $P_A, P_B, \rho,
\rho_{AR}$ of the disturbance process. These parameters are not known
\emph{a priori} but emerge from operational measurement. Hence:
\begin{equation}
  \pi^* \;\notin\; \{f \,:\, f \text{ is closed-form analytic in
  observable signals only}\}.
\end{equation}
\end{theorem}

\textbf{Interpretation.} An offline tuner cannot derive $\pi^*$ by
analysis alone; an online procedure that fits a model to data is
required. We do not claim this rules out alternative learning
paradigms (e.g., system-identification, Kalman filters with adaptive
covariance) but we do claim that ``hand-tuning'' is provably
insufficient.

\subsection{Theorem 8: AI Necessity (Main Result)}

We synthesise the preceding into our main theoretical result.

\begin{theorem}[AI Necessity Theorem]
\label{thm:ai-necessity}
Suppose the following operationally realistic assumptions hold:
\begin{itemize}\setlength\itemsep{1pt}
\item \textbf{(A1) Non-stationary traffic}: $P_t \in \{P_A, P_B\}$ with
mode-switch rate $\rho > 0$ (Theorem~\ref{thm:static-regret}).
\item \textbf{(A2) Moment-matched Byzantine}: at least one disturbance
realisation matches the first two moments of legitimate traffic
(Theorem~\ref{thm:suff-stat}).
\item \textbf{(A3) Measurement autocorrelation}: legitimate
measurements carry $\mathrm{AR}(1)$ autocorrelation $\rho_{AR} \neq 0$
(Theorem~\ref{thm:memory}).
\end{itemize}
Then:
\begin{enumerate}
\item No constant policy $\pi \in \Pi_{const} \cup \Pi_{lin}^{mem-0}$
achieves bounded regret w.r.t.\ $\pi^*$
(Theorem~\ref{thm:static-regret}, \ref{thm:memory}).
\item No linear windowed policy $\pi \in \Pi_{lin}^{mem-K}$ achieves
positive Byzantine-detection capacity
(Theorem~\ref{thm:suff-stat}).
\item No closed-form analytic policy expresses $\pi^*$
(Theorem~\ref{thm:intractability}).
\end{enumerate}
Consequently, the unique known function class that admits
sample-complexity-bounded convergence to $\pi^*$ is data-fitted
non-linear windowed functions, instantiated by neural networks.
\end{theorem}

\textbf{This is the central thesis of the paper.} Machine learning is
not an optimisation: it is the \emph{logically unique known route} to
deadline-miss-bounded regret under (A1)--(A3). Sections~\ref{sec:predictor}
and~\ref{sec:safety} then implement and bound the chosen ML instance.

\subsection{Quantitative Decomposition of Deadline-Miss Rate}

The deadline-miss rate $\eta(\pi)$ decomposes as
\begin{equation}
  \eta(\pi) = \eta_{cascade}(\pi) + \eta_{byz}(\pi) + \eta_{degrade}(\pi),
\end{equation}
where each term is bounded below for hand-tuned policies by Theorems
\ref{thm:static-regret}--\ref{thm:memory} respectively:
\begin{itemize}\setlength\itemsep{1pt}
\item $\eta_{cascade}$ is lower-bounded by the static-regret term of
Theorem~\ref{thm:static-regret}, scaling with $\epsilon^2 \rho$.
\item $\eta_{byz}$ is lower-bounded by the constant gap implied by
Theorem~\ref{thm:suff-stat}'s zero-information capacity.
\item $\eta_{degrade}$ is lower-bounded by Theorem~\ref{thm:memory}'s
memoryless-detector deficit.
\end{itemize}
ML drives all three terms to operational levels simultaneously
(Sections~\ref{sec:empirical-ceiling}, \ref{sec:exp-results}); no
hand-tuned policy can do all three.

\subsection{Why \emph{Transformers}, Specifically}

Theorem~\ref{thm:ai-necessity} requires non-linear windowed
modelling but does not single out Transformers. Among ML
architectures, attention is preferred for one specific reason.

\begin{proposition}[Cross-Channel--Cross-Time Representation Efficiency]
\label{prop:attn-eff}
For input $\mathbf{x} \in \mathbb{R}^{K \times d}$, expressing all
interactions of order up to $r$ between any subset of channels and any
subset of timesteps requires parameter counts:
$\Omega((Kd)^r)$ for a flat MLP, $\Omega(d^r s^{r-1})$ for CNNs
with kernel size $s$, $\Omega(d^r)$ for RNNs (with $O(K)$ sequential
depth), and $\Omega(d^r)$ for self-attention with $O(\log K)$ parallel
depth.
\end{proposition}

At $K{=}60$ and $d{=}8$, self-attention reaches the most favourable
parameters-per-expressivity ratio. This is why our 141\,k-parameter
Transformer reaches AUC$=0.927$ (Table~\ref{tab:ceiling-emp}) at
$0.56$\,MB.

\subsection{Defence Against Theoretical Objections}

\textbf{``OGD on $\Pi_{lin}^{mem-K}$ adapts online.''}
Online gradient descent navigates \emph{within} the policy class.
Theorem~\ref{thm:suff-stat} bounds the class itself, so online
adaptation cannot exceed the zero-information ceiling.

\textbf{``Tree-based methods are non-linear; do we need attention?''}
We tested a 100-tree random forest in Section~\ref{sec:ablation},
obtaining AUC$=0.806$---still 12\,pp below the Transformer. Trees
satisfy Theorem~\ref{thm:ai-necessity} (they are non-linear, windowed,
data-fit) but at higher parameter cost. Theorem 8 selects ML as a
class; Proposition~\ref{prop:attn-eff} selects Transformers within ML.

\textbf{``Are assumptions (A1)--(A3) too strong?''}
(A1) is empirically observed in every multi-AZ AWS production trace we
have access to. (A2) is the standard Kerckhoffs adversary in
security---assuming the attacker has read the paper. (A3) is measured
($\rho_{AR} \in [0.4, 0.7]$) in network and OS-scheduling traces.
All three are conservative, not strong.

\textbf{``Is the closed-form bound of Theorem~\ref{thm:intractability}
provable?''}
We claim the weak form: no rational function in observable parameters
expresses $\pi^*$. Stronger forms involving algebraic or analytic
function classes require deeper theory; we leave the precise
algebraic-complexity bound to future work.
